Navegava certa vez por uma rede social qualquer quando tropecei no relato sobre um professor que orientava seus alunos a lrem como se mirassem o mar e, ocasionalmente, avistassem baleias emergindo e saltando através de sua superfície ondulante;
    sem qualquer preocupação em enxergar todos os seus mistérios, mas atentos àqueles com os quais se deparassem:
        orientação que reconheço como uma das mais necessárias para a leitura de textos e mundos que venho realizando desde então.

Mais recentemente, li uma carta de Bruno Latour -- falecido pensador francês das técnicas e ciências com o qual me associo na tarefa de \textit{composição de um \textbf{mundo comum}} -- endereçada a uma indefinida \textit{senhorita} -- aparentemente uma aluna, quicá uma orientanda -- inquieta e confusa quanto ao curso do professor Latour;
    ele tenta apaziguá-la elaborando os assuntos, a abordagem e a proposta do curso, mas admite ser necessário "\textit{ir até o fim para compreender do que se trata}":
        um jeito certamente curioso de apaziguar o coração de uma estudante cheia de dúvidas no qual reconheci, entretanto, uma das orientações mais necessárias para me acompanhar na escrita que inicio.

Ainda mais recentemente, li um relato de Givânia Maria da Silva -- quilombola educadora do quilombo de Conceição das Crioulas -- sobre a história transmitida oralmente de seu território;
    ela conta das mulheres que chegaram ali no ano de 1802 e compraram aquelas terras com o fruto de seu trabalho fiando algodão para vender na cidade, e que um homem chegou depois trazendo uma imagem de Nossa Senhora da Conceição na intenção de convertê-las ao Catolicismo, mas que "elas trouxeram a santa para o lado delas, houve um pacto feminino com a santa, e o homem perde a importância naquela história, ninguém sabe mais o que aconteceu":
        reconheço aqui uma sabedoria necessária, feminista e poderosa para aprender como ficar com o problema.

O problema com o qual escolhi ficar na pesquisa que relato e realizo neste texto é o \textbf{objetivismo};
    este é o nome que proponho para articular o interesse pessoal que me remeteu a este trabalho com o interesse coletivo necessário à realização de uma pesquisa científica:
        ao menos segundo a ética e gramática científicas que almejo propor.

Minha jornada teve início no dia dois de abril de 2024;
    até então, vinha tentando ingressar num programa de mestrado para fazer uma leitura do romance "Viva o povo brasileiro", de João Ubaldo Ribeiro, que discutisse o papel das construções simbólicas e históricas para a configuração social do Brasil contemporâneo:
        contudo, naquela data -- Dia Mundial de Conscientização do Autismo -- a publicação de uma página de divulgação científica criticando a expressão pública de pessoas autistas acerca de nossas experiências e reduzindo o diagnóstico a a uma confluência com a ordem neoliberal despertou em mim a necessidade de encontrar uma maneira de falar sobre as contingências do autismo sem cair nos reducionismos característicos da prática científica moderna.

Meu objetivo era -- e continua sendo -- conseguir especular e descrever um mundo em que não seja inadequado ser \textit{assim};
    um objetivo bastante pessoal para uma pesquisa científica, exatamente o que Latour propõe em um dos critérios de seu "Princípio de Falseabilidade de Stengers-Despret" para guiar uma ciência propositiva e capaz de fomentar a composição de um mundo comum:
        quanto mais próximo e interessado estiver o investigador científico de seu problema de pesquisa, mais e melhor ele lhe oferecerá escuta.

Mas eu preciso de uma baleia para o público não autista se espero que ele acompanhe a leitura até o final para compreender do que estou tratando;
    para falar, eu preciso ser ouvido, e para ser ouvido, eu preciso trazer o público para o meu lado, e para isso proponho um pacto contro o objetivismo:
        um \textbf{pacto contra-objetivista}.

Neste pacto, caberá a você que me lê prosseguir com atenção não só às baleias, mas também às controvérsias, contradições, paradoxos, ambiguidades, incompletudes e incoerências;
    quanto a mim, fico comprometido a conduzir-lhes por \textbf{uma jornada em busca de uma nova bússola para comunicações -- e práticas -- científicas bem articuladas}:
        associo assim ao sucesso de meu objetivo a proposição de uma gramática científica -- como diria Nego Bispo -- \textbf{politeísta}.

Se por um lado faço o melhor que posso nesta introdução para apresentar um ponto de partida inteligível, por outro carrego comigo o mesmo problema de Latour com seu curso;
    deixo desde já advertido que este é um texto espiralar, não linear: